% ─── Capítulo 6: Código Intermedio ────────────────────────────────────────────
\chapter{BACKEND — Código Intermedio (TAC)}

\section{Fundamentos Teóricos}

Una vez que el AST ha sido optimizado, el compilador necesita traducirlo a una forma más lineal que pueda ser procesada por fases sucesivas o por distintos backends. Esta forma intermedia es el \textbf{Código de Tres Direcciones (TAC, \textit{Three-Address Code})}.

Teóricamente, el TAC modela un tipo de máquina abstracta llamada \textbf{Máquina de Registros Infinitos}: tiene una cantidad ilimitada de variables temporales (\texttt{t0}, \texttt{t1}, \texttt{t2}...) que actúan como registros de cómputo. Esto permite \textit{linearizar} cualquier expresión anidada del AST sin tener que preocuparse aún por la limitación de registros físicos de la CPU real.

Cada instrucción TAC puede expresarse formalmente como una \textbf{cuádrupla}:
\[ (\mathrm{op},\, \mathrm{arg_1},\, \mathrm{arg_2},\, \mathrm{result}) \]
Por ejemplo:
\begin{itemize}
    \item \texttt{t2 = t0 + t1} → $(\texttt{+},\; \texttt{t0},\; \texttt{t1},\; \texttt{t2})$
    \item \texttt{if\_false t2 goto L1} → $(\texttt{if\_false},\; \texttt{t2},\; \texttt{L1},\; -)$
    \item \texttt{call duplicar} → $(\texttt{call},\; \texttt{duplicar},\; -,\; -)$
\end{itemize}

\begin{infobox}[¿Por qué es necesario el Código Intermedio?]
El TAC es un nivel de abstracción fundamental porque:
\begin{enumerate}
    \item \textbf{Portabilidad:} Al ser independiente de la arquitectura (x86, ARM, RISC-V), el mismo TAC puede ser reutilizado para generar código para distintos procesadores simplemente intercambiando el backend.
    \item \textbf{Diagnóstico:} Permite al programador o al docente leer y comprender qué operaciones realiza el compilador antes de sumergirse en el ensamblador.
    \item \textbf{Base para optimizaciones posteriores:} Técnicas avanzadas como la Asignación Global de Registros (GRA), el Análisis de Flujo de Datos y la Eliminación de Subexpresiones Comunes (CSE) operan sobre formas de TAC o SSA.
\end{enumerate}
\end{infobox}

\begin{notebox}
En el compilador TPP, el código intermedio tiene propósito \textbf{diagnóstico y educativo}. No produce un archivo de salida propio; el TAC se imprime en pantalla para que el programador pueda entender cómo el compilador «ve» el programa antes de generar el ensamblador final.
\end{notebox}

\section{Implementación: \texttt{ir.c}}

\subsection{Formato TAC}

Cada instrucción TAC sigue uno de estos formatos:

\begin{verbatim}
t0 = 10               // Asignacion de constante a temporal
t1 = t0 + t2          // Operacion binaria
x = t1                // Asignacion a variable
arr[t1] = t0          // Escritura en arreglo
t0 = arr[t1]          // Lectura de arreglo
t0 = x < 5            // Comparacion
if t0 goto L1         // Salto condicional
goto L2               // Salto incondicional
L1:                   // Etiqueta
call nombre_fun       // Llamada a función
ret                   // Retorno
print "texto"         // Impresión de cadena
print t0              // Impresión de variable
\end{verbatim}

\subsection{Generador Recursivo}

\begin{lstlisting}[style=cstyle, caption={Generación de TAC recursiva en ir.c}]
static int contador_temp = 0;
static int contador_etiq = 0;

void generar_tac(NodoAST *nodo) {
    if (!nodo) return;

    switch (nodo->tipo) {
        case N_ENTERO:
            printf("    t%d = %d\n", contador_temp++, nodo->val_entero);
            break;

        case N_VARIABLE:
            printf("    t%d = %s\n", contador_temp++, nodo->nombre_var);
            break;

        case N_OPERACION: {
            generar_tac(nodo->izq);
            int temp_izq = contador_temp - 1;
            generar_tac(nodo->der);
            int temp_der = contador_temp - 1;
            const char *op = obtener_nombre_operador(nodo->operador);
            printf("    t%d = t%d %s t%d\n", 
                   contador_temp++, temp_izq, op, temp_der);
            break;
        }

        case N_MIENTRAS: {
            int etiq_inicio = contador_etiq++;
            int etiq_fin = contador_etiq++;
            printf("L%d:\n", etiq_inicio);
            generar_tac(nodo->izq);  // Condicion
            printf("    t%d = %d\n", contador_temp, contador_temp);
            printf("    if_false t%d goto L%d\n", contador_temp-1, etiq_fin);
            generar_tac(nodo->der);  // Cuerpo
            printf("    goto L%d\n", etiq_inicio);
            printf("L%d:\n", etiq_fin);
            break;
        }

        case N_DECLARACION_FUN:
            printf("fun %s:\n", nodo->nombre_var);
            generar_tac(nodo->der);   // Cuerpo
            printf("end_fun %s\n", nodo->nombre_var);
            break;

        case N_LLAMADA_FUNCION:
            printf("call %s\n", nodo->nombre_var);
            printf("    t%d = t-1\n", contador_temp++); // Valor de retorno
            break;

        case N_RETORNAR:
            generar_tac(nodo->izq);
            printf("ret\n");
            break;
    }
}
\end{lstlisting}

\section{Ejemplo de TAC}

Para el programa:
\begin{lstlisting}[style=tppstyle]
entero x = 5;
entero y = x + 3;
imprimir(y);
\end{lstlisting}

El TAC generado sería:
\begin{verbatim}
    t0 = 5
    x = t0
    t1 = x
    t2 = 3
    t3 = t1 + t2
    y = t3
    print y
\end{verbatim}

Para una función:
\begin{verbatim}
fun duplicar:
    t0 = x
    t1 = 2
    t2 = t0 * t1
    res = t2
ret
end_fun duplicar
\end{verbatim}

\section{TAC para Estructuras de Control}

\subsection{Condicional si/sino}

\begin{verbatim}
    t0 = x
    t1 = 0
    t2 = t0 > t1         // t2 = condicion
    if_false t2 goto L1  // Si falso, ir a L1 (sino)
    // bloque si
    goto L2
L1:
    // bloque sino
L2:
\end{verbatim}

\subsection{Bucle para}

\begin{verbatim}
    t0 = 0
    i = t0               // inicializacion: i = 0
L1:
    t1 = i
    t2 = 5
    t3 = t1 < t2         // condicion: i < 5
    if_false t3 goto L2
    // cuerpo del bucle
    t4 = i
    t5 = 1
    t6 = t4 + t5
    i = t6               // paso: i = i + 1
    goto L1
L2:
\end{verbatim}
